% ---
% id: reductions/rem:why-exact
% kind: rem
% title: "Exactness of the reductions and generality of the transfer"
% summary: >-
%   Equality of full outputs lets one argument transfer every criterion with
%   the same ε, while a distance bound would need a modulus of continuity no
%   criterion is given. Every reduction and link of this contract, the
%   domination link included, is exact.
% prominence: standard
% status: accepted
% uses: [def:opreduction, lem:opreduction-transfer, thm:transfer, def:criterion, thm:factorization, thm:chain, thm:emit]
% mentions: [rem:emitclaim, lem:source-transfer, meat-core/prop:optol-nobound, transmitter/def:source-segment, req:no-approximate-relations, req:mixed-status-chain, cor:reference-companion, reductions/domination/thm:noise-channel, reductions/domination/cor:dom-chain]
% assumes: []
% refs: []
% provenance: newly-derived
% lean:
% notes: "meat-core/prop:optol-nobound bounds the gap it exhibits by its own representability condition (N at most 52); do not describe that gap as arbitrary."
% ---

\begin{remark}\label{rem:why-exact}
Definition~\ref{def:opreduction} demands equality of the \emph{full} output
state, and that demand is what pays for everything
Lemma~\ref{lem:opreduction-transfer} delivers. A security criterion in the
sense of Definition~\ref{def:criterion} is an arbitrary function of the output
state: no linearity, no monotonicity and no continuity is assumed of it
anywhere in this contract. Under exact equality that costs nothing, because
part (i) of the lemma is $f$ applied to two equal states; under anything
weaker it costs everything, because a hypothesis
$\|\mathsf{out}(a)-\mathsf{out}'(r(a))\|_1\le\delta$ bounds
$f(\mathsf{out}(a))$ in terms of $f(\mathsf{out}'(r(a)))$ only through a
modulus of continuity for $f$, which the definition of a criterion does not
give and which the consuming engine's criterion would have to supply case by
case. So the generality of the transfer, one argument for every criterion with
the same $\varepsilon$, and the exactness are one fact, not two.

\emph{Both instances are exact, for a checkable reason.} On the receiver side
the substitution rests on the channel identity \eqref{eq:factorization}, an
identity of maps with no error term (Theorem~\ref{thm:factorization}). On the
source side it rests on the source identity (S3) of
\xref{transmitter/def:source-segment}, an exact matrix identity on declared
numbers: it is decided entry by entry where every \tref{source link} supplies
its channel as data, and concluded from a named claim where a link supplies
its channel through a \tref{theorem constructor}
(Lemma~\ref{lem:source-transfer}).

\emph{What an approximate reduction would need.} It would need a theorem
converting a distance between output states into a security penalty, and this
contract has none. The security transfer, Theorem~\ref{thm:transfer}, rests on
an equality and has no penalty slot, so there is no parameter for such a term
to enter; accordingly a declared relation carrying an error term is refused
(Requirement~\ref{req:no-approximate-relations}), and a \tref{reduction chain}
mixing an exact link with one carrying an error term is refused rather than
given a composite status, a status for the chain as a whole
(Theorem~\ref{thm:chain}\,(iv), Requirement~\ref{req:mixed-status-chain}).
The nearest thing to a continuity statement in these contracts is a refusal of
one: \xref{meat-core/prop:optol-nobound} exhibits two
\tref{exactly admissible} target families within a declared distance of each
other whose rate bounds per round differ by an amount limited only by the
number model. The limit is that proposition's own representability condition,
which binary64 arithmetic imposes and which caps the integer $N$ measuring the
gap at $N\le52$; so proximity of models licenses nothing about rates there
either.

\emph{The \tref{emission displacement} is not an error term of a reduction.}
The displacement $\Delta_s$ of the \tref{emission normalizer}, the largest
change it makes to an entry of an instance's datum in the \tref{data sector}
$s$, does not make a reduction approximate. It moves the \tref{declared model}
the reduction is stated about, and the reduction's identity then holds of that
model with no error term at all (Theorem~\ref{thm:emit},
Remark~\ref{rem:emitclaim}). Nor does the step back to the
\tref{handed model}, the construction before the normalizer's \tref{snap},
need an error term. No bound relates a key length computed from the declared
model alone to one computed from the handed model, and the displacement is
not what relates them (Theorem~\ref{thm:emit}\,(v)). The gap is not bridged
at all where a run can avoid it: a run that forms the number its key length
is formed from on the reference companion
(Corollary~\ref{cor:reference-companion}), built with no normalizer from the
\tref{handed reference}, or from the instance's own snapped blocks where none
is handed, reads the declared model only to find points. Where a run does read it, the gap is closed by a disclosed
premise about the device, never by an error term. The \tref{domination link},
which no requirement emits, is exact as well: a noise channel followed by its
target with classical flags (Theorem~\ref{thm:noise-channel}), so that a
squashing link followed by it is exact as a whole
(Corollary~\ref{cor:dom-chain}).
\end{remark}
